%\documentclass[wcp,gray]{jmlr} % test grayscale version
\documentclass[wcp]{jmlr}


% givasile packages
\usepackage{bbm}
\usepackage{xfrac}

% givasile commands
\newcommand{\xc}{\mathbf{x}_c}
\newcommand{\Xc}{\mathbf{X}_c}
\newcommand{\X}{\mathbf{X}}
\newcommand{\Ysloc}{Y_{s}^{\text{local}}}
\newcommand{\ysloc}{y_{s}^{\text{local}}}
\newcommand{\x}{\mathbf{x}}
\newcommand{\xii}{\mathbf{x}^{(i)}}
\newcommand{\xsi}{x_s^{(i)}}
\newcommand{\xci}{\mathbf{x}_c^{(i)}}
\newcommand{\yi}{y^{(i)}}
\newcommand{\E}{\mathbb{E}}
\newcommand{\var}{\mathrm{Var}}
\newcommand{\R}{\mathbb{R}}
\newcommand{\1}{\mathbbm{1}}
\newcommand{\D}{\mathcal{D}}

% The following packages will be automatically loaded:
% amsmath, amssymb, natbib, graphicx, url, algorithm2e

%\usepackage{rotating}% for sideways figures and tables
\usepackage{longtable}% for long tables

% The booktabs package is used by this sample document
% (it provides \toprule, \midrule and \bottomrule).
% Remove the next line if you don't require it.
\usepackage{booktabs}
% The siunitx package is used by this sample document
% to align numbers in a column by their decimal point.
% Remove the next line if you don't require it.
%\usepackage[load-configurations=version-1]{siunitx} % newer version
%\usepackage{siunitx}
%\usepackage{natbib}

% Do not comment the following commands:
\pagenumbering{gobble}
\newcommand{\cs}[1]{\texttt{\char`\\#1}}
\makeatletter
\let\Ginclude@graphics\@org@Ginclude@graphics 
\makeatother

\jmlrvolume{}
\jmlryear{2022}
\jmlrworkshop{ACML 2022}

\title[Short Title]{Full Title of Article}

 % Use \Name{Author Name} to specify the name.
 % If the surname contains spaces, enclose the surname
 % in braces, e.g. \Name{John {Smith Jones}} similarly
 % if the name has a "von" part, e.g \Name{Jane {de Winter}}.
 % If the first letter in the forenames is a diacritic
 % enclose the diacritic in braces, e.g. \Name{{\'E}louise Smith}

 % Two authors with the same address
 % \author{\Name{Author Name1} \Email{abc@sample.com}\and
 %  \Name{Author Name2} \Email{xyz@sample.com}\\
 %  \addr Address}

 % Three or more authors with the same address:
 % \author{\Name{Author Name1} \Email{an1@sample.com}\\
 %  \Name{Author Name2} \Email{an2@sample.com}\\
 %  \Name{Author Name3} \Email{an3@sample.com}\\
 %  \Name{Author Name4} \Email{an4@sample.com}\\
 %  \Name{Author Name5} \Email{an5@sample.com}\\
 %  \Name{Author Name6} \Email{an6@sample.com}\\
 %  \Name{Author Name7} \Email{an7@sample.com}\\
 %  \Name{Author Name8} \Email{an8@sample.com}\\
 %  \Name{Author Name9} \Email{an9@sample.com}\\
 %  \Name{Author Name10} \Email{an10@sample.com}\\
 %  \Name{Author Name11} \Email{an11@sample.com}\\
 %  \Name{Author Name12} \Email{an12@sample.com}\\
 %  \Name{Author Name13} \Email{an13@sample.com}\\
 %  \Name{Author Name14} \Email{an14@sample.com}\\
 %  \addr Address}


 % Authors with different addresses:
  \author{\Name{Author Name1} \Email{abc@sample.com}\\
  \addr Address 1
  \AND
  \Name{Author Name2} \Email{xyz@sample.com}\\
  \addr Address 2
 }

\editors{Emtiyaz Khan and Mehmet Gonen}

\begin{document}

\maketitle

\begin{abstract}
This is the abstract for this article.
\end{abstract}
\begin{keywords}
List of keywords separated by semicolon.
\end{keywords}

\section{Introduction}
Main contents here.

The decision about the bin size has major influence in the feature
effect plot. Therefore, it is very important (a) to inform the user to
what extend they should trust the ALE plot and, consequently, decide
the bin size with the optimal effect. In the end, we will see that it i

\newpage
~\newpage

\section{Related Work}

\newpage

\section{Probabilistic formulation of ALE plots}

\subsection{ALE background}\label{sec:ALE}

This section introduces the reader to the ALE formulation of feature
effect. Given that \(f: \R^S \rightarrow \R \) is known, we can
measure the effect of the \(s\)-th feature at a specific point
\( \x = (\xc, x_s)\) of the input space \(\mathcal{X}\) as
\(f_s(\x) = \sfrac{\partial f(\x)}{\partial x_s}\). ALE models the
\textit{local} effect at \(x_s\) as the expected effect over the
distribution of the unkown (latent) features \(\Xc\), i.e.
\(\E_{p(\xc;x_s)}[f_s(x_s, \xc)]\). Afterwards, the \textit{global}
effect at \(x_s\) is the accumulation of the \textit{local} effects:
%
\begin{equation}
  \label{eq:ale-definition}
  f_{\mathtt{ALE}}(x_s) =
  \int_{x_{s, \text{min}}}^{x_s} \E_{p(\xc;x_s=z)}[f_s(x_s, \xc)] \partial z
\end{equation}

In real cases, it is infeasible to compute
eq.~\eqref{eq:ale-definition} analytically. Therefore, we reside on
estimating the effect from the training set. Let's denote the
available dataset as \(\D = \{\xii, \yi \}_{i=1}^N\), where \( \xii \)
is the \(i\)-th feature vector of the training data and \(\yi\) is the
\(i\)-th label.~\cite{Apley2020} proposed spliting the axis into
\( K \) equal-sized bins, find the set of points that lie in each bin,
i.e.  \( \mathcal{S}_k = \{ \xii : x_s^{(i)} \in [z_{k-1}, z_k) \} \)
and, finally, find the local effect at each bin as the mean value of
the population \(\hat{\mu}_{s,k}\). The global effect is then
estimated through the following formula:
%
\begin{equation}
  \label{eq:ale-approximation}
  \hat{f}_{\mathtt{ALE}}(x_s)
  = \Delta x \sum_{k=1}^{k_{x_s}} \hat{\mu}_{s,k}
  = \sum_{k=1}^{k_{x_s}} \frac{1}{|\mathcal{S}_k|} \sum_{i:\xsi \in \mathcal{S}_k} [f(z_k, \xci) - f(z_{k-1}, \xci)]
\end{equation}
%
where \(\Delta x\) is the bin size and \(k_{x_s}\) the index of the
bin that includes \(x_s\).
%

\subsection{Quantification of the uncertainty of the estimation}

We propose an alternative estimation of the feature effect inside each
bin. The alternative estimation is better because it (a) separates the
decision about the bin limits from the local effect of each point, (b)
secures from out-of-distribution sampling and (c) provides better
variance estimation inside the bin. The alternative estimation is:

\begin{equation}
  \label{eq:1}
  \tilde{f}_{\texttt{ALE}}(x_s) =
  \Delta x \sum_{k=1}^{k_{x_s}} \tilde{\mu}_{s,k} =
  \Delta x \sum_{k=1}^{k_{x_s}} \frac{1}{|\mathcal{S}_k|} \sum_{i:\xsi \in \mathcal{S}_k} f_s(\xii)
\end{equation}
%
Our goal is to inform the user to what extent they should trust the
ALE plot. For this purpose, we introduce two metrics; the variance and
the standard error of the estimation. The variance of the estimation
is:

\begin{equation}
  \label{eq:effect-var-approx}
  \tilde{\sigma}_s^2(x_s)
  = (\Delta x)^2 \sum_{k=1}^{k_{x_s}} \tilde{\sigma}_{s,k}^2
  = (\Delta x)^2 \sum_{k=1}^{k_{x_s}} \frac{1}{|\mathcal{S}_k|} \sum_{i:\xii \in \mathcal{S}_k} (f_s(\xii) - \tilde{\mu}_{s,k})^2 
\end{equation}
%
The standard error of the estimation is:

\begin{equation}
  \label{eq:effect-std-error}
  \text{SE}(x_s) = \Delta x \sum_{k=1}^{k_{x_s}} (\frac{\hat{\sigma}_{s,k}^2}{|\mathcal{S}_k|})^\frac{1}{2}
\end{equation}


The two metrics, (\ref{eq:effect-var-approx}) and
(\ref{eq:effect-std-error}), should be trusted only in case they
respect some contstraints. Under these constraints, the user can be
confident that the correct feature effect is inside two or three
standard errors errors around the mean estimation.


Our computations are based on the hypothesis that inside
all bins \textbf{the gradient wrt. to the feature of interest doesn't
  depend on the value feature of interest}. Unfortunately, we cannot
when this is the case. We just know that as the bins grow larger, it
is more possible to violated this hypothesis. In this case both the
ALE effect and the standard error are wrong.

The standard error should be trusted when it is estimated by a large
of data points. For example, in plots (c) and (d), there are bins with
less than 10 points. Therefore, in these cases, we cannot trust the
plot or the standard error.

\subsection{Unsupervised metric for assessing the quality of ALE plots}

Based on the observations of the previous chapter, we want to create a
single metric for choosing the best feature effect plot. Our goal is
to minimize the variance, given that we have points inside eac bin 

We propose the minimization of the accumulated standard deviation (or
accumulated variance). For ALE with K bins, let's notate as:

\begin{itemize}
\item \(dx^K\) the length of each bin
\item \(p_i^K \) the number of the training points inside the \(i-th\) bin
\item \(\sigma_i^K \) the std of the local effects of the training points inside the \(i-th\) bin
\end{itemize}
%
We will minimize:

\begin{align} \label{eq:min-criterion}
  K_{min} &= \text{argmin}_{K} \quad  [dx^K  \sum_i^K \sigma_i^K * (1 - d_i^K)]\\
  & \text{s.t. } p_i^K \geq \texttt{min\_points\_per\_bin } \forall i
\end{align}
%
where \(d_i^K = 0.2 * \frac{p_i^K}{N} \in [0,0.1] \) works as a `discount', favoring the
creation of bigger bins in cases of similar standard deviation.

\subsection{ALE plots with variable-size bins}

\newpage
~\newpage

\section{Experiments}

\subsection{Synthetic Data sets}

\begin{enumerate}
\item One example to show that (a) variance and (b) standard error are
  good estimates to what extend to trust the ALE plot.
\item One example to show the importance of variable size bins
\end{enumerate}

\newpage
~\newpage

\subsection{Real Data sets}

\newpage

\section{Conclusion}


A figure in Fig.~\ref{fig:spiral}. Please use high quality graphics
for your camera-ready submission -- if you can use a vector graphics
format such as \texttt{.eps} or \texttt{.pdf}.
\begin{figure}[htp]
\begin{center}
\includegraphics[width=0.5\textwidth]{spiral.eps}
\caption{A spiral.}\label{fig:spiral}
\end{center}
\end{figure}

An example of citation~\cite{DBLP:conf/acml/2009}.

\newpage

%\acks{Acknowledgements should go at the end, before appendices and references. You can uncomment this for the camera-ready version on paper acceptance.}

%\bibliographystyle{plain}
\bibliography{acml22}
\newpage

\appendix

\section{Symbols}\label{apd:symbols}

Spaces and points:

\begin{itemize}
\item \(\mathcal{X} \subseteq \R^S \): the input space
\item \(\x = (x_s, \xc) \in \mathcal{X}\)
\item \(x_s \in \R \)
\item \(\xc \in \R^{S-1}\)
\item
\end{itemize}
%
Functions:
\begin{itemize}
\item \(f: \R^S \rightarrow \R\): black-box function
\item \(f_s(\x) = \sfrac{\partial f(\x)}{\partial x_s}: \R^S \rightarrow \R\): effect at point \(\x\)
\item 
\end{itemize}
%
ALE:
\begin{itemize}
\item \(\sfrac{f(z_{k_x}) - f(z_{k_x-1})}{\Delta x}\): effect of \(x_s\) at \(\x\)
\item \(\E_{p(\xc;x_s)}[f_s(x_s, \xc)]\): local effect of \(x_s\) at \(x_s\)
\item \(f_{\texttt{ALE}}(x_s): \R \rightarrow \R\): global effect of \(x_s\) at \(x_s\)
\item \(\hat{f}_{\texttt{ALE}}(x_s): \R \rightarrow \R\): \textit{estimator} of global effect of \(x_s\) at \(x_s\)
\end{itemize}

\section{Derivations}\label{apd:derivations}

Derivations

\begin{multline}
    \Ysloc \sim p(\ysloc ; x_s)
    = \int{} p(\ysloc|\xc;x_s) p(\xc;x_s) \partial \xc \\
    = \int_{\xc} \delta(\ysloc-f_s(\x)) p(\xc;x_s) \partial \xc
    = \int_{\xc} \1(\ysloc=f_s(\x)) p(\xc;x_s) \partial \xc
\end{multline}

Some statistics for the local feature effect:

\begin{equation}
  \E[ \Ysloc ; x_s] = \int_{\xc} p(\xc|x_s) f_s(\x) \partial \xc
\end{equation}

\begin{equation}
  \var[ \Ysloc ; x_s] = \int_{\xc} p(\xc|x_s) (f_s(\x) - \E[ \Ysloc ; x_s])^2 \partial \xc
\end{equation}

Some statistics for the global feature effect:

\begin{equation}
  \E[Y_s ; x_s] = \int_{x_{s,min}}^{x_s} \E(\ysloc;z) \partial z
\end{equation}

\begin{equation}
  \var[Y_s ; x_s] = \int_{x_{s,min}}^{x_s} \var[ \Ysloc ; x_s] (\partial z)^2
\end{equation}


\section{Second Appendix}\label{apd:second}

This is the second appendix.


\end{document}



% \subsection{Probabilistic formulation of ALE plots}

% In this section, we will introduce a probabilistic formulation for ALE
% plots. As discussed before, ALE directly defines the local effect as
% the expectated value over the unknown features,
% i.e. \(\E_{p(\xc;x_s)}[f_s(x_s, \xc)]\). Therefore, the local and the
% global effect, which is the integration of the local effects, are
% modeled as deterministic variables.

% For adjusting ALE plots into a probabilistic framework, we model the
% local effect with a random variable conditioned on feature \(x_s\),
% i.e., \(\Ysloc \sim p(\ysloc ; x_s)\):

% \begin{equation}
%     p(\ysloc ; x_s)
%     = \int_{\xc} \delta(\ysloc-f_s(\x)) p(\xc;x_s) \partial \xc
%     \label{eq:local-effect-rv}
% \end{equation}
% %
% Following the core idea of ALE, we define the global feature effect as
% the accumulation of the local effects. Therefore, the feature effect
% of the \(s\)-th feature is \(Y_s \sim p(y_s ; x_s) \):
% \begin{equation}
%   p(y_s ; x_s) = \int_{x_{s,min}}^{x_s} p(\ysloc;x_s=z) \partial z
%   \label{eq:global-effect-rv}
% \end{equation}
% %
% In eqs.~(\ref{eq:local-effect-rv}), (\ref{eq:global-effect-rv}) we
% reformulate ALE plots in a fully probabilistic manner. We are
% particularly interested into two summary statistics of \(Y_s\), the
% expected value \(\mu_s (x_s) = \E[Y_s ; x_s]\) and the variance
% \(\sigma_s^2(x_s) = \var[Y_s ; x_s]\). In the appendix, we show that
% the expected value coincides with the definition of ALE plots in
% eq.~(\ref{eq:ale-definition}), i.e.
% \(\mu_s (x_s) = f_{\texttt{ALE}(x)s}\). The variance can be computed
% as:

% \begin{equation}
%   \sigma_s^2 (x_s)
%   = \int_{x_{s,min}}^{x_s} \E_{p(\xc|x_s = z)}[(f_s(\x) - \mu_s(z))^2](\partial z)^2
% \end{equation}


% As before, it is infeasible to compute the the variance
% \(\sigma_s^2(x_s)\) in analytically. Therefore, we reside on
% estimating them from the available examples of the training set:

% \begin{equation}
%   \label{eq:effect-var-approx}
%   \hat{\sigma}_s^2(x_s)
%   = (\Delta x)^2 \sum_{k=1}^{k_{x_s}} \hat{\sigma}_{s,k}^2
%   = (\Delta x)^2 \sum_{k=1}^{k_{x_s}} \frac{1}{|\mathcal{S}_k|} \sum_{i:\xii \in \mathcal{S}_k} (f_s(\xii) - \hat{\mu}_{s,k})^2 
% \end{equation}

% .. add a description for standard error

% \begin{equation}
%   \label{eq:effect-std-error}
%   \text{SE}(x_s) = \Delta x \sum_{k=1}^{k_{x_s}} (\frac{\hat{\sigma}_{s,k}^2}{|\mathcal{S}_k|})^\frac{1}{2}
% \end{equation}
